% !TeX root = ../colloquium-2.tex

% Формулировки теорем Ролля и Лагранжа.

\phantomsection\label{sec_Definition12}
\begin{definition}[Формулировки теорем Ролля и Лагранжа.]
    \subsection*{Теорема Ролля.}
    \begin{note}[purple]
        Пусть:
        \begin{enumerate}
            \item Функция $f$ определена и непрерывна на отрезке $[a,\; b]$
            \item Функция $f$ дифференцируема на интервале $(a, \; b)$
            \item $f(a) = f(b)$
        \end{enumerate}
        Тогда существует такая точка $c \in (a, \; b)$, что $f'(c) = 0$. 
    \end{note}

    \subsection*{Теорема Лагранжа.}
    \begin{note}[purple]
        Пусть:
        \begin{enumerate}
            \item Функция $f$ определена и непрерывна на отрезке $[a,\; b]$
            \item Функция $f$ дифференцируема на интервале $(a, \; b)$
        \end{enumerate}
        Тогда существует такая точка $c \in (a, \; b)$, что $f(b) - f(a) = f'(c)(b - a)$.
    \end{note}
\end{definition}